% Part: normal-modal-logic
% Chapter: syntax-and-semantics
% Section: relational-models

\documentclass[../../../include/open-logic-section]{subfiles}

\begin{document}

\olfileid{nml}{syn}{rel}

\olsection{النماذج العلائقية}

مدار دلالة المنطقات الجهية النظامية على \emph{النموذج العلائقي}.
نأخذ مجموعة من العوالم، ونجعل بينها علاقة ثنائية تسمى علاقة الوصول،
ثم نعيّن أي !!{propositional variable}s تصدق في كل عالم. فبهذه
الأجزاء الثلاثة يتحدد النموذج.

\begin{defn}
  \emph{نموذج} اللغة الموجهية الأساسية ثلاثية~$\mModel{M}
  = \tuple{W, R, V}$، حيث
  \begin{enumerate}
  \item $W$ مجموعة غير خالية من ``العوالم''،
  \item $R$ علاقة وصول ثنائية على~$W$، و
  \item $V$ دالة تعين لكل !!{propositional variable}~$p$ مجموعة
    $V(p)$ من العوالم الممكنة.
  \end{enumerate}
  وعندما يصدق~$Rww'$، نقول إن~$w'$ \emph{يمكن الوصول إليه من}~$w$.
  وعندما يكون~$w \in V(p)$، نقول إن~$p$ \emph{صادق في}~$w$.
\end{defn}

ومن فائدة الدلالة العلائقية أن نموذجها يُرى في رسم يسير، كما
في~\olref{fig:simple}. فالعوالم عقد، ووصول العالم~$w'$ من~$w$
يمثله سهم يخرج من~$w$ وينتهي إلى~$w'$. ونضع على العقدة
$\mTrue{p}$ حين يكون~$w \in V(p)$، ونضع~\mFalse{p} فيما عدا ذلك.
وعليه فإن~\olref{fig:simple} يرسم النموذج الذي تعيّنه البيانات
$W = \{w_1, w_2, w_3\}$، و$R = \{\tuple{w_1, w_2},\tuple{w_1,
  w_3}\}$، و$V(p) = \{w_1, w_2\}$، و$V(q) = \{w_2\}$.

\begin{figure}
  \begin{center}
    \begin{tikzpicture}[modal]
      \node[world] (w1) [label={[align=right]right:\mTrue{p}\\ \mFalse{q}}]{$w_1$};
      \node[world] (w2) [label={[align=right]right:\mTrue{p}\\ \mTrue{q}},
        above right=of w1]{$w_2$};
      \node[world] (w3) [label={[align=right]right:\mFalse{p}\\ \mFalse{q}},
        below right=of w1] {$w_3$};
      \draw[->] (w1) to (w2);
      \draw[->] (w1) to (w3);
    \end{tikzpicture}
  \end{center}
  \caption{نموذج بسيط.}
  \ollabel{fig:simple}
\end{figure}

\end{document}
